%\pagebreak
\cvsection{主要獲獎與榮譽}

NAACL 2022 最佳論文獎（最佳新任務）暨最佳論文獎（最佳新資源） \hfill \textcolor{black}{2022}\\ %-2016\\
% NAACL 2022 最佳論文獎（最佳新資源） \hfill \textcolor{black}{2022}\\ %-2016\\
NAACL 2022 傑出論文獎（以人為本自然語言處理特別主題貢獻獎） \hfill \textcolor{black}{2022}\\ %-2016\\
AMTA 2022 最佳簡報獎 \hfill \textcolor{black}{2022}\\ %-2016\\
% 史丹佛以人為本人工智慧研究補助（研究專案 HabitLab） \hfill \textcolor{black}{2018}\\ %-2016\\
% 麻省理工學院網頁程式設計競賽（6.470）決賽入圍暨佳作，2013\\ % (Project: PsetParty)
%ACM UIST（使用者介面軟體與技術）學生創新競賽「最實用獎」第一名 \hfill \textcolor{sectcol}{2012}\\
% ACM UIST（使用者介面軟體與技術）學生創新競賽第一名 \hfill \textcolor{black}{2012}\\
% ACM CHI（人機互動）學生研究競賽第一名 \hfill \textcolor{black}{2012}\\
%波士頓音樂駭客日 InstantKaraoke 專案第一名 \hfill \textcolor{sectcol}{2012}\\ % (Project: InstantKaraoke)
美國國家科學基金會研究生研究獎學金（NSF Graduate Fellowship） \hfill \textcolor{black}{2013}\\ %(declined in favor of NDSEG), 2013\\
美國國防科學與工程研究生獎學金（NDSEG Graduate Fellowship） \hfill \textcolor{black}{2013}\\ %-2016\\
Phi Beta Kappa（麻省理工學院前 10\% 學生）與 Tau Beta Pi（麻省理工學院工程學院前 12.5\% 學生）榮譽會員 \hfill \textcolor{black}{2012}\\
% Phi Beta Kappa 學會會員（麻省理工學院前 10\% 學生） \hfill \textcolor{black}{2012}\\
% Tau Beta Pi 學會會員（麻省理工學院工程學院前 12.5\% 學生）、Eta Kappa Nu 學會會員（麻省理工學院電機資訊學院前 25\% 學生） \hfill \textcolor{black}{2011}\\
% 麻省理工學院自主機器人競賽（MASLAB）第一名 \hfill \textcolor{black}{2010}\\
